\section{Day 3: Minimal Polynomials (Sep.\ 15, 2026)}
Office hours will be held from $3$--$4$pm on Wednesdays in BA6110. Last time, we defined the algebraic numbers, $\QQ^{\opname{alg}}$, and the algebraic integers, $\ZZalg$. % yeah we're going to use these macros from now on </3
We also saw that $\QQalg$ is a field, and $\ZZalg$ is a ring.

\begin{proposition}
    Let $\alpha \in \QQalg$ be nonzero; there is a unique monic irreducible polynomial $f(x) \in \QQ[x]$ such that $f(\alpha) = 0$. Moreover, if $g(x) \in \QQ[x]$ such that $g(\alpha) = 0$, then $f \mid g$.
\end{proposition}
\begin{definition} \label{defn:minimal_polynomial}
    Let $\alpha \in \QQalg$. Then the unique monic irreducible $f(x) \in \QQ[x]$ with $f(\alpha) = 0$ is called the \textit{minimal polynomial of $\alpha$}. The \textit{degree of $\alpha$} is the degree of $f(x)$ itself.
\end{definition}

If $\alpha \in \QQalg$, write $\QQ(\alpha) = \{h(\alpha) \mid h \in \QQ[x]\}$. If $\alpha \in \ZZalg$, write $\ZZ[\alpha] = \{h(\alpha) \mid h \in \ZZ[x]\}$. Clearly, $\QQ(\alpha) \subset \CC$, and we call it the subring generated by $\QQ$ and $\alpha$ in $\CC$. Similarly, $\ZZ(\alpha)$ is the smallest subring of $\CC$ containing $\alpha$.

As an example, take $\QQ(\sqrt{-5}) = \{a + b \sqrt{-5} \mid a, b \in \QQ\}$, and $\ZZ[\sqrt{-5}] = \{a + b\sqrt{-5} \mid a, b \in \ZZ\}$.

\begin{proposition} \label{prop:minimal_polynomial_properties}
    Let $\alpha \in \QQalg$.
    \begin{enumerate}[(i)]
        \item $\QQ(\alpha) = \QQ[x]/p(x)$, where $p$ is the minimal polynomial of $\alpha$.
        \item As a vector space over $\QQ$, $\QQ(\alpha)$ has basis $1, \alpha, \dots, \alpha^{n-1}$, where $n = \deg(\alpha)$.
        \item $\QQ(\alpha)$ is a field.
    \end{enumerate}
\end{proposition}
\begin{proof}
    Consider the map $\ev_\alpha \colon \QQ[x] \to \QQ(\alpha)$, given by $h \mapsto h(\alpha)$. $\ev_\alpha$ is a surjective ring homomorphism by definition of $\QQ(\alpha)$ (cf.\ earlier), whence
    \[ \QQ(\alpha) \cong \frac{\QQ[x]}{\ker(\ev_\alpha)}. \]
    In particular, we saw that $\ker(\ev_\alpha) = \{h \in \QQ[x] \mid h(x) = 0\} = \left<f(x)\right>$, where $f(x) = a_0 + \dots + a_{n-1} x^{n-1} + x^n$, so we are done with (i). (ii) immediately follows from earlier. For (iii), pick $\beta \in \QQ(\alpha)$ nonzero, and let $\beta = h(\alpha)$ with $h \in \QQ[x]$, for which $f \nmid h$, since $h(\alpha) \neq 0$. By the division algorithm (or Bezout's lemma), there must exist polynomials $A, B \in \QQ[x]$ such that $Af + Bh = 1$. Evaluating at $\alpha$, we see that $A(\alpha)f(\alpha) = 0$, so $B(\alpha)h(\alpha) = 1$, so we have found an inverse $\beta\inv = B(\alpha) \in \QQ(\alpha)$ as desired.
\end{proof}

\begin{example}
    We give some examples of irreducible polynomials.
    \begin{enumerate}[(i)]
        \item The minimal polynomial of $\sqrt{-5}$ is $x^2 + 5$.\footnote{so... any idea why this is irreducible? ``the eisenstein criterion?'' *grimaces* wayyyy overkill...}
        \item Let $\omega = \frac{1}{2} (-1 + \sqrt{-3}) = \exp(2\pi i / 3)$. Clearly, $\omega^3 = 1$. Since $x^3 - 1$ is not irreducible, we see that we may factor and obtain the actual minimal polynomial for $\omega$ to be $x^2 + x + 1$.
    \end{enumerate}
\end{example}

Is $\frac{1}{2} \in \ZZalg$? Recall Gauss' lemma; suppose $R$ is a UFD. A polynomial $f(x) \in \RR[x]$ is \textit{primitive} if the GCD of its coefficients is $1$; Gauss' lemma states that if $f, g \in R[x]$ is primitive, then $fg$ is primitive as well. As a quick proof, suppose not; then there exists $p \in R$ prime such that $p \mid fg$; since $R/p$ is an integral domain, let $\ol f, \ol g \in (R/p)[x]$ be the corresponding reductions. We have that $\ol f \cdot \ol g = 0$, but then it must be that $\ol f = 0$ or $\ol g = 0$, which is a contradiction.

\newpage
\begin{lemma} \label{lem:coefficients_in_ring}
    Let $R$ be a UFD, with $F$ its field of fractions. Suppose $g, h \in F[x]$ are monic such that $f \coloneq gh \in R[x]$. Then $g, h \in R[x]$.
\end{lemma}
\begin{proof}
    Let $c, d$ be the LCM of the denominators of the coefficients of $g, h$. Then $cg$ and $dh$ are primitive polynomials in $R[x]$. Gauss' lemma would then imply that $cg \cdot dh = (cd) \cdot gh$ is a primtiive polynomial; but $gh \in R[x]$, so $cd \mid (cd) \cdot gh$ in $R[x]$, i.e., $cd$ must be a unit in $R$. It follows that $c$ and $d$ are both units in $R$, whence $g, h \in R[x]$.
\end{proof}

\begin{corollary} \label{cor:minpoly_of_alg_int_is_in_Z[x]}
    If $\alpha \in \ZZalg$, then its minimal polynomial lies in $\ZZ[x]$.
\end{corollary}
\begin{proof}
    Let $p(x) \in \QQ[x]$ be the minimal polynomial of $\alpha$. Since $\alpha \in \ZZalg$, there exists a monic $h(x) \in \ZZ[x]$ such that $h(\alpha) = 0$. $p(x) \mid h(x)$, i.e., $h(x) = p(x) q(x)$, for some $q$ with rational coefficients. From here, the lemma tells us that $p, q \in \ZZ[x]$.
\end{proof}

\begin{corollary} \label{cor:intersection_of_Q_and_alg_int}
    $\QQ \cap \ZZalg = \ZZ$.
\end{corollary}
\begin{proof}
    Let $\alpha \in \QQ \cap \ZZalg$. The minimal polynomial of $\alpha$ is $x - \alpha$, so $\alpha \in \ZZ$.
\end{proof}

\begin{definition} \label{defn:conjugate}
    Algebraic numbers $\alpha, \beta \in \QQalg$ are \textit{conjugate} if they have the same minimal polynomial.
\end{definition}

\begin{lemma} \label{lem:conjugate_roots}
    We have some more properties of minimal polynomials.
    \begin{enumerate}[(i)]
        \item If $f \in \QQ[x]$ is monic and irreducible of degree $d$, then $f$ has $d$ distinct roots in $\CC$.
        \item If $\alpha$ is an algebraic number of degree $d$, then there exist exactly $d$ algebraic numbers conjugate to $\alpha$.
    \end{enumerate}
\end{lemma}
Clearly, (i) implies (ii). We will prove this next lecture.